%  06_implementation.tex  –  Chapter 5: Implementation

\chapter{Implementation}
\label{ch:implementation}

\section{Raspberry Pi Subsystem}
\label{sec:rpi-subsystem}

\subsection{Startup and Initialization}
\label{subsec:startup}

When the station script is launched, a series of initialization steps take place before data
collection begins. The SPI interface is opened and configured for communication with the MCP3208
ADC at a maximum clock speed of 1{,}350{,}000\,Hz. DHT11 initialization takes place on GPIO
pin 23. GPIO pins of the two LEDs and the button are initialized using the \texttt{gpiozero}
library. A hardware interrupt is attached to GPIO pin 24 of the KY-037 digital output, which
operates in both rising and falling edge modes with 50\,ms debounce period to ensure that the
same sound is not counted twice because of bouncing.

Then the script makes two external API calls. First to \texttt{ipinfo.io} so it can determine
the device's current city and country based on public IP address. Second one to prompt the user
to enter a neighbourhood label through the terminal. The variable neighbourhood label holds a
string value, which is incorporated into each data record until the end of the session. Lastly,
the principal thread for collecting data is initiated as a daemon thread, after which the script
blocks to wait for input from the button.

\vspace{12pt}
\begin{figure}[H]
  \centering
  \includegraphics[width=0.6\textwidth]{startup_flowchart}
  \caption{Figure 13: Flowchart of the station startup sequence}
  \label{fig:startup-flowchart}
\end{figure}
\vspace{12pt}

\subsection{Data Reading}
\label{subsec:data-reading}

Every 30 minutes, it is checked whether data sending is enabled via the thread-safe event flag.
If it is, the loop proceeds to read all sensors in sequence.

First we read the DHT11 sensor since the temperature and humidity data obtained from it will be
used as inputs for calculating the PPM values from the MQ sensors using the correction equation.
As the communication of DHT11 follows the single-wire protocol, which is time-sensitive and
often does not produce any output, the read is repeated up to five times with an interval of two
seconds. In case of failure of all the attempts, then the process is simply skipped. It has been
found that the DHT11 fails to provide any reading when continuously operated for around 14 hours.

The five MQ sensors and the KY-037 analog output are then read through the MCP3208 ADC. The
MCP3208 communicates over SPI and returns a 12-bit integer value between 0 and 4095 for each
channel. This raw value is converted to a voltage using the reference voltage of 3.3\,V:

\vspace{12pt}
\begin{figure}[H]
  \centering
  \includegraphics[width=0.75\textwidth]{raw_voltage_code}
  \caption{Figure 14 -- Converted raw voltage code}
  \label{fig:raw-voltage-code}
\end{figure}
\vspace{12pt}

However, because the MQ sensors operate on 5\,V and their analog output passes through a
1\,k$\Omega$--2\,k$\Omega$ voltage divider before reaching the ADC, the voltage read by the
ADC is not the actual sensor output voltage. The actual sensor output must be recovered by
applying the inverse of the divider ratio:

\vspace{12pt}
\begin{figure}[H]
  \centering
  \includegraphics[width=0.75\textwidth]{actual_voltage_code}
  \caption{Figure 15 -- Actual output value code}
  \label{fig:actual-voltage-code}
\end{figure}
\vspace{12pt}

This corrected voltage is then used for all subsequent PPM calculations.

\vspace{12pt}
\begin{figure}[H]
  \centering
  \includegraphics[width=0.65\textwidth]{voltage_divider}
  \caption{Figure 16: Diagram of the voltage divider circuit for a single MQ sensor, showing the 1\,k$\Omega$ and 2\,k$\Omega$ resistors, the 5\,V sensor output and the scaled voltage reaching the MCP3208 input.}
  \label{fig:voltage-divider}
\end{figure}
\vspace{12pt}

\subsubsection{PPM Computation}
\label{subsubsec:ppm}

MQ sensors work by making their internal resistance vary ($R_S$) in response to changes in gas
concentration. The relation between resistance and concentration is described by a power-law
curve published in the sensor datasheets:

\begin{equation}
  \mathrm{PPM} = a \cdot \left(\frac{R_S}{R_O}\right)^b \tag{1}
  \label{eq:ppm}
\end{equation}

$a$ and $b$ are sensor-specific constants. They are derived from the datasheet sensitivity
curves. $R_O$ is the sensor's resistance measured in clean air, that serves as the calibration
baseline. $R_S$ is computed from the measured output voltage:

\begin{equation}
  R_S = \frac{(V_C - V_{out}) \cdot R_L}{V_{out}} \tag{2}
  \label{eq:rs}
\end{equation}

$V_C = 5\,\text{V}$ is the sensor supply voltage and $R_L = 2000\,\Omega$ is the load
resistance.

The $R_O$ values used for this project were derived from an initial set of data that was
collected under mostly clean indoor air conditions. The first baseline was recorded during the
early phase of the testing; however, it was discovered later on that the $R_O$ values were
inaccurate because of inadequate sensor pre-heating and were later re-calibrated. The details
specific to this show in section 6.1.1 (Table 7). In both cases, the same steps were taken to
find $R_S$ from the average measured voltage using Eqs.\ (2) and (3), where $R_O = R_S$.

\begin{equation}
  V_{out} = V_{ADC} \cdot \frac{R_1 + R_2}{R_2} \tag{3}
  \label{eq:vdivider}
\end{equation}

The recalibrated $R_O$ values and constants for each sensor can be seen in the following table:

\vspace{12pt}
\begin{table}[H]
  \centering
  \caption{Table 2: Constant values for gases}
  \label{tab:sensor-constants}
  \renewcommand{\arraystretch}{1.3}
  \begin{tabular}{lllrr}
    \toprule
    \textbf{Sensor} & \textbf{Target gas} & \textbf{$a$} & \textbf{$b$} & \textbf{$R_O$ ($\Omega$)} \\
    \midrule
    MQ135 & eCO\textsubscript{2}  & 110.47 & $-2.862$ & 30{,}688 \\
    MQ3   & Alcohol/Benzene        & 0.3934  & $-1.504$ &  7{,}524 \\
    MQ6   & LPG/Butane             & 1000.5  & $-2.180$ &  8{,}490 \\
    MQ7   & Carbon Monoxide        & 99.042  & $-1.518$ & 11{,}539 \\
    MQ8   & Hydrogen               & 976.97  & $-0.688$ & 38{,}502 \\
    \bottomrule
  \end{tabular}
\end{table}
\vspace{12pt}

MQ sensor resistance is also affected by temperature and humidity, so a correction factor which
results from Hanwei application notes~\cite{hanwei} is applied to $R_S$:

\begin{equation}
  \text{CF}(T,RH) = \frac{1}{(-0.01223 + 0.00609 \cdot RH) \cdot T + (1.70836 - 0.14289 \cdot RH)} \tag{4}
  \label{eq:correction}
\end{equation}

$T$ is temperature in Celsius and $RH$ is relative humidity percentage. 20$^\circ$C and 65\,\%
RH are the references. The corrected resistance is then used in the ratio:

\vspace{12pt}
\begin{figure}[H]
  \centering
  \includegraphics[width=0.75\textwidth]{ppm_code}
  \caption{Figure 17 -- PPM computation code}
  \label{fig:ppm-code}
\end{figure}
\vspace{12pt}

Results outside the range of 0 to 100{,}000 PPM are discarded as invalid and stored as null
values.

\vspace{12pt}
\begin{figure}[H]
  \centering
  \includegraphics[width=0.65\textwidth]{ppm_flowchart}
  \caption{Figure 18: Flowchart of the PPM computation pipeline}
  \label{fig:ppm-flowchart}
\end{figure}
\vspace{12pt}

\subsubsection{Sound Detection}

The KY-037 sound sensor provides two outputs. The analog output (AOUT) is read through the
MCP3208 on channel 5 and converted to a voltage, representing the background ambient sound
level at the time of the reading. The digital output (DOUT) is connected directly to GPIO
pin 24 and monitored via a hardware interrupt registered with \texttt{RPi.GPIO}. The interrupt
works by triggering during both the rising and falling edges of the signal, incrementing the
global counter for every trigger. The trigger counter value is saved and reset to zero every
30 minutes. Interrupts are used instead of software polling, as they can detect very brief
signal transitions that occur within microseconds, which software polling can not catch due to
its slow processing speed.

\subsubsection{Warning Conditions and LED Indicators}

At the end of each measurement cycle, the computed PPM values are compared against a set of
approximate safety thresholds sourced from publicly available environmental guidelines:

\vspace{12pt}
\begin{table}[H]
  \centering
  \caption{Table 3 -- PPM threshold for each sensor}
  \label{tab:ppm-thresholds}
  \renewcommand{\arraystretch}{1.3}
  \begin{tabular}{ll}
    \toprule
    \textbf{Sensor} & \textbf{Threshold} \\
    \midrule
    MQ135 & 1{,}200 PPM \\
    MQ3   & 50 PPM      \\
    MQ6   & 1{,}600 PPM \\
    MQ7   & 35 PPM      \\
    MQ8   & 1{,}000 PPM \\
    \bottomrule
  \end{tabular}
\end{table}
\vspace{12pt}

In the case where at least one sensor exceeds its threshold, a warning message is generated
naming the sensors that triggered it. Additionally, weather-based warnings are generated if the
current weather condition is stormy or unknown, if the temperature exceeds 35$^\circ$C under
sunny conditions, or is below 0$^\circ$C. When a warning is active, the red LED is turned on
for 5 seconds. The blue LED is on whenever data is successfully sent.

\subsubsection{Data Packaging, Local Logging, and Transmission}

After each cycle, all data collected and calculated are put together in a dictionary object
with Python consisting of 30 keys: timestamp divided into year, month, and day keys, location
keys, temperature in Celsius and Fahrenheit, humidity, weather description, raw ADC data,
calculated voltages, PPM data, sound data, and a warning key if there is any. The dictionary
object gets saved locally on a CSV file by using \texttt{csv.DictWriter} function in Python.
This way a local copy of all recorded data is created, so in case something goes wrong while
sending data to the server, there will always be the local file. The dictionary also gets
transformed into JSON format and is sent through an HTTP POST request to the Flask
application's endpoint \texttt{/api/sensor}.

\vspace{12pt}
\begin{figure}[H]
  \centering
  \includegraphics[width=0.65\textwidth]{data_cycle_flowchart}
  \caption{Figure 19: Flowchart of a single complete data collection cycle}
  \label{fig:data-cycle}
\end{figure}
\vspace{12pt}

\subsection{Prediction Script}
\label{subsec:prediction-script}

The prediction part has two components: a training executed in Google Colab and a daily
inference script (\texttt{predict.py}) deployed on the Raspberry Pi.

\subsubsection{Training Data and Feature Engineering}

The machine learning models were trained in Google Colab by using three data sources. These
were combined into a single training dataset:

Weather data for Timișoara covering January 2024 to the present was fetched from Open-Meteo
archive API~\cite{openmeteo}. This provided temperature and humidity.

Official air quality values were downloaded from the EEA Air Quality Download
Service~\cite{eea}. They cover four government monitoring stations in Timișoara (RO0193A,
RO0194A, RO0196A, RO0197A). The data goes over 17 months with hourly data and includes CO,
NO\textsubscript{2}, SO\textsubscript{2} and O\textsubscript{3}. Only valid observations were
converted to daily data by averaging of all the stations.

The live values measured by the sensors from the Raspberry Pi sensor station (MQ135, MQ3, MQ6,
MQ7 and MQ8 PPM values, temperature and humidity) gave us the most accurate data which we used
to predict the level of risk.

The input daily data was extended with the use of lagged variables (yesterday's data and last
week's data), average of the last seven days for temperature, humidity, and gas concentrations
and time variables such as day of the year, month, season, and week of the year.

\subsubsection{Model 1 — Prophet (Temperature and Humidity Forecast)}

Meta's Prophet library~\cite{prophet} was used to train two separate time series models. One
for temperature and one for humidity, each producing a next 7-day forecast. Prophet was chosen
for its ability to handle missing data gaps, multiple seasonality patterns (weekly and yearly)
and suitability for datasets with limited history (perfect for this type of project). One model
was trained per target variable and exported as a \texttt{.pkl} file using \texttt{joblib}.

\subsubsection{Model 2 — Random Forest Classifier (Danger Prediction)}

Random Forest classifier~\cite{sklearn} was trained to predict if a future day is likely to
present dangerous air quality conditions. The danger label was defined as any MQ sensor reading
going above its established safety threshold. Where live sensor data was available, the label
was derived directly from the sensor readings. For days without sensor coverage, the label was
derived from the EEA pollutant readings compared against WHO, EU, and OSHA guideline
thresholds. The classifier was trained with balanced class weights to account for the natural
imbalance between safe and dangerous days in the dataset and evaluated using 5-fold
cross-validation for estimating generalization performance.

\subsubsection{Deployment and Retraining}

Both trained models are exported as lightweight \texttt{.pkl} files and placed in a
\texttt{models/} subdirectory on the Raspberry Pi alongside a \texttt{model\_meta.json} file
containing the list of feature column names expected at inference time.

The \texttt{predict.py} file is called on a daily basis at 7:00 AM via a cron job. On each
run the script imports both models, gets a current 7-day weather forecast from the Open-Meteo
forecast API~\cite{openmeteo}, creates features vectors for each of the following 7 days
(based on the local sensor data from a CSV file), generates danger predictions and forecasts,
writes the predictions into a local \texttt{predictions.csv} file. It then POSTs predictions
to the Flask server on AWS and stores them in the \texttt{aq\_predictions} table in PostgreSQL.

The system can be retrained monthly (or even more often) by re-executing one notebook in Google
Colab with updated sensor data and historical data from EEA and Open-Meteo, in order to
improve model performance.

The steps of each daily inference run are logged to a \texttt{predict.log} file in the station
directory. This way a record is created which also includes any errors that appear.

\vspace{12pt}
\begin{figure}[H]
  \centering
  \includegraphics[width=0.65\textwidth]{predict_flowchart}
  \caption{Figure 20 -- \texttt{predict.py} daily inference pipeline}
  \label{fig:predict-flowchart}
\end{figure}
\vspace{12pt}

\subsection{ML Training Notebook}
\label{subsec:ml-notebook}

The ML models are trained in a specialized Google Colab notebook~\cite{googlecolab}
(\texttt{AirQuality\_ML\_Training.ipynb}). This notebook contains 11 stages. It should be run
from top to bottom, without interruption, in one go. Workflow consists of: data collection
(three different sources), data preparation, model training (and evaluation), making predictions
for the next 7 days and exporting the models. A special cell for model retraining is provided
for the monthly update of models based on new data from the sensors CSV file.

\subsubsection{Historical Weather Data}

Timișoara related historical weather data is fetched automatically from the Open-Meteo archive
API~\cite{openmeteo} using the coordinates (45.7489$^\circ$N, 21.2087$^\circ$E), covering
January 2021 to the current date. The API returns hourly temperature, relative humidity, and
WMO weather codes, which are resampled to daily averages. The most frequent weather code of
each day is retained as the daily code, and a binary \texttt{is\_rainy} flag is derived by
mapping all WMO codes of 51 or above to rainy days.

\subsubsection{EEA Station Data}

The CSV files for the EEA data set are uploaded via the Colab file uploader tool. The pollutant
type is automatically determined using the numeric pollutant code present in each of the
filenames as follows: 00001 for SO\textsubscript{2}, 00007 for O\textsubscript{3}, 00010 for
CO, and 00020 for NO\textsubscript{2}. All readings whose validity flag is 1 and have positive
and non-sentinel values are used.

\subsubsection{Sensor CSV}

The sensor CSV file exported from the Raspberry Pi station is uploaded through the Colab file
picker. The date is parsed from the station's format combining the Year, Month, and the numeric
portion of the Day field. The remaining rows are aggregated to daily mean and maximum PPM values
for each of the five MQ sensors, alongside daily mean temperature and humidity. The danger label
for sensor-covered days is derived by checking whether any MQ sensor's daily mean PPM exceeds
its threshold:

\vspace{12pt}
\begin{table}[H]
  \centering
  \caption{Table 4 -- Sensor threshold value and its basis}
  \label{tab:ml-thresholds}
  \renewcommand{\arraystretch}{1.3}
  \begin{tabular}{lll}
    \toprule
    \textbf{Sensor} & \textbf{Threshold} & \textbf{Basis} \\
    \midrule
    MQ135 & 200 PPM   & EU NO\textsubscript{2} hourly limit (proxy) \\
    MQ3   & 1{,}000 PPM & OSHA TWA \\
    MQ6   & 1{,}000 PPM & $\approx$20\,\% lower explosive limit \\
    MQ7   & 35 PPM    & WHO 1-hour guideline \\
    MQ8   & 1{,}000 PPM & $\approx$25\,\% lower explosive limit \\
    \bottomrule
  \end{tabular}
\end{table}
\vspace{12pt}

\subsubsection{Merging and Feature Engineering}

The three daily datasets are merged, with the weather data as the base. Short gaps of up to
three consecutive days in the EEA data are forward-filled to account for brief station outages.
The merged dataset is then enriched with the following features:

Time-based features: day of year, month, week of year, and season (winter, spring, summer,
autumn mapped from month number).

Lag features: yesterday's and last week's values for temperature, humidity, and each EEA
pollutant, giving the model awareness of short and medium-term trends.

Rolling average: rolling mean over 7 days for temperature, humidity, and each EEA pollutant to
smooth the data and capture weekly patterns.

\subsubsection{Exploratory Plots}

The script creates six time series plots in PNG format. Each plot shows the daily pattern for
temperature, humidity, NO\textsubscript{2}, CO, SO\textsubscript{2} and O\textsubscript{3}
during the entire date range. This is used for sanity checking the data set after merging.

\subsubsection{Prophet Training}

Two Prophet models are trained. One for temperature and one for humidity. It is used the
\texttt{train\_prophet()} function. The settings for each model include yearly and weekly data
in additive form with a \texttt{changepoint\_prior\_scale} value of 0.05 to make the trend
more flexible. The last 14 days of each series are held out to use as a test dataset. MAE
(Mean Absolute Error) and RMSE (Root Mean Squared Error) are calculated for the dataset and
printed. Each model provides an extended forecast for 7 days ahead of the training dataset.
Prophet forecast plots are saved as PNG files showing the fitted values, the 7-day projection,
and the uncertainty interval.

\subsubsection{Random Forest Training}

The danger label is built with sensor-derived values taking priority over EEA-derived values,
where both are available. The final feature set for the classifier includes all weather
features, time features, lag and rolling features, EEA pollutant features, and MQ sensor PPM
means, where available. Rows missing the danger label are dropped. Remaining missing values in
the feature columns are filled with the column median before training.

The Random Forest classifier is configured with 200 estimators, maximum depth of 6, minimum 3
samples per leaf, balanced class weights and a fixed random seed of 42. It is evaluated using
5-fold stratified cross-validation with F1 score as the metric. The cross-validation F1 mean
and standard deviation are printed. The classifier is then fitted on the full dataset for
deployment. Feature importances are computed and the top 15 are plotted and saved as a PNG.

% Flask Server 

\section{Flask Server}
\label{sec:flask-server}

\subsubsection*{Sensor Data Ingestion}

The \texttt{/api/sensor} endpoint receives POST requests from the Raspberry Pi and inserts the
data into the \texttt{sensor\_data} table in the PostgreSQL database. All numeric fields are
passed through a \texttt{safe\_float()} helper that returns \texttt{None} instead of raising an
exception if a value cannot be converted (for cases where Raspberry Pi sends the string
``N/A'' when sensors return invalid readings). Database connections are managed through a
\texttt{SimpleConnectionPool} provided by \texttt{psycopg2}. It maintains between 1 and 10
open connections and reuses them across requests to avoid the overhead of opening a new
connection for each incoming reading.

\subsubsection*{User Authentication}

The website supports user registration, email verification, login, logout and password reset.
Passwords are not stored in plain text. All of them are hashed using bcrypt before being saved
to the database. The function used for verification is \texttt{checkpw} (from Bcrypt).

\vspace{12pt}
\begin{figure}[H]
  \centering
  \includegraphics[width=0.5\textwidth]{encrypted_password}
  \caption{Figure 21 -- How the encrypted password looks in DB}
  \label{fig:encrypted-pwd}
\end{figure}
\vspace{12pt}

Upon a user registering, a secure, random token is generated by using Python's
\texttt{secrets} module. It is then stored with the unverified account. A verification email
is sent to the provided address via Gmail SMTP using STARTTLS encryption. The email contains a
link to the \texttt{/api/verify-email/<token>} endpoint. When the link is clicked, the token
is looked up in the database, the account marked as verified and the token cleared. Login is
blocked for accounts that haven't been verified.

Password reset follows the same token-based pattern: reset token is generated, stored in a
separate table with a one-hour expiry timestamp, then sent to the user by email. When the reset
form is submitted, the token is validated against the database and verified for expiry, before
the new password is accepted and hashed.

A hardcoded administrator account also exists. It is checked before the database lookup at
login time. The admin session grants access to the admin panel. This panel displays registered
user activity, the form data and allows toggling of a maintenance mode for the entire website.
When maintenance mode is active, non-admin users see a maintenance notice instead.

\vspace{12pt}
\begin{figure}[H]
  \centering
  \includegraphics[width=0.65\textwidth]{registration_flowchart}
  \caption{Figure 22 - Flowchart of the user registration and email verification flow}
  \label{fig:registration}
\end{figure}
\vspace{12pt}

\subsubsection*{Sensor Data Querying and Filtering}

The \texttt{/api/sensor-data} endpoint accepts optional query parameters for year, month, day,
neighbourhood, city and country. They are used to dynamically build a parameterized WHERE
clause in the SQL query. Only parameters that are present in the request are included in the
filter. Omitting all parameters returns the most recent unfiltered 500 records. The
\texttt{/api/sensor-filters} endpoint returns the distinct values present in the database for
each filterable field, which the frontend uses to populate the filter dropdowns.

\subsubsection*{CSV Export}

The \texttt{/api/download-csv} endpoint requires an active user session, returning a 401 error
if the request comes from an unauthenticated user. When authenticated, it queries all rows from
the \texttt{sensor\_data} table ordered by insertion time, writes them to an in-memory
\texttt{StringIO} buffer. Python's \texttt{csv} module is used for it. After that it returns
the buffer contents as a file download.

\subsubsection*{Predictions API}

Two endpoints handle the prediction data. The POST endpoint \texttt{/api/predictions} receives
daily forecast from the Raspberry Pi. It needs a secret to be shared via the
\texttt{X-Secret} header, where the server then validates the value against a constant before
accepting any data. This way unauthorized writing to the predictions table is prevented. The
endpoint receives a body in JSON format containing an array of prediction objects, with each
object inserted in the \texttt{aq\_predictions} table in an upsert operation using
\texttt{ON CONFLICT (date) DO UPDATE}. This means that if the station executes the prediction
script on the same day multiple times, only the first insertion would be done.

The GET endpoint \texttt{/api/predictions} is publicly available and returns 7 forecast objects
with dates starting from the current date. This ensures that the website displays the next 7
days' predictions, not the previous ones.

\subsubsection*{Admin Messages Endpoint}

\texttt{/api/admin/messages} API call has the \texttt{require\_admin} decorator and retrieves
all rows in the \texttt{contact\_messages} table. It has an optional order query parameter that
determines whether the results should be sorted in ascending or descending order based on the
time of submission, and its default value is \texttt{desc}. The returned fields include the
message id, sender name, email, subject, body, and creation date.

%  Website Frontend 

\section{Website Frontend}
\label{sec:frontend}

\subsection{Overall Structure}
\label{subsec:frontend-structure}

The website is composed of multiple pages, each served as a static HTML file by the Flask
server. All pages share a navigation header. Each page has its own inline JavaScript
responsible for client-side logic.

\subsubsection*{Home Page}

The home page introduces the project, lists the hardware components, and provides an overview
of the system features. It is rendered from \texttt{index.html} and contains only static
content and images.

\vspace{12pt}
\begin{figure}[H]
  \centering
  \includegraphics[width=0.85\textwidth]{homepage}
  \caption{Figure 23: Home page}
  \label{fig:homepage}
\end{figure}
\vspace{12pt}

\subsubsection*{Data Page}

The Data page's inline JavaScript handles filter population, data fetching and chart rendering
entirely on the client side.

On page load, two asynchronous requests are made in sequence. The first calls
\texttt{/api/sensor-filters}, which returns the distinct values present in the database for
year, month, day, neighbourhood, city and country. These values are used to populate six
dropdown menus organized into two filter rows. There is one for time-filtering and another for
location filtering. Entries with blank or ``blank'' values are ignored. The second API call
retrieves the unfiltered dataset from \texttt{/api/sensor-data} containing maximum 500 entries
sorted by the time of their insertion.

When a user applies filter values and presses the Apply button, the \texttt{loadData()}
function generates a query string based on selected dropdowns and retrieves the filtered
dataset. When a user clicks Reset All button, all dropdowns will be cleared and the unfiltered
dataset will be retrieved.

Chart drawing is done using Chart.js 4.4.1~\cite{chartjs}, fetched from a CDN. Eight line
charts are drawn, each chart representing one measured parameter. The construction of these
charts is done by the function \texttt{buildChart()}. This makes sure that all the instances
of Chart.js are destroyed prior to creation to avoid memory leakage caused by multiple
instances.

Each dataset in a chart is constructed using the function \texttt{makeDataset()}. It is
included in the script of \texttt{data.html}:

\vspace{12pt}
\begin{figure}[H]
  \centering
  \includegraphics[width=0.85\textwidth]{makedataset_code}
  \caption{Figure 24: makeDataset() code}
  \label{fig:makedataset}
\end{figure}
\vspace{12pt}

As there are more than 100 points, they are not shown by setting \texttt{pointRadius = 0}.
This helps to make the chart readable for larger sets of data. The value of the tension of
0.35 gives some curvature to the line. The null values in the data set, for the invalid
sensors that were recorded in the database as None, are passed as \texttt{null} to the
JavaScript arrays and handled by the library as gaps in the line.

The sound chart is a special case, as here are drawn two datasets on two separate Y axes. The
left Y axis shows the KY-037 analog voltage in volts and the right Y axis shows the number of
events.

Each chart has their visual style defined in the \texttt{CHART\_DEFAULTS} and \texttt{COLORS}
constants at the top of the script.

\vspace{12pt}
\begin{figure}[H]
  \centering
  \includegraphics[width=0.85\textwidth]{data_charts}
  \caption{Figure 25: Data page charts}
  \label{fig:data-charts}
\end{figure}
\vspace{12pt}

\begin{figure}[H]
  \centering
  \includegraphics[width=0.6\textwidth]{data_flow}
  \caption{Figure 26: Data Page workflow}
  \label{fig:data-flow}
\end{figure}
\vspace{12pt}

\subsubsection*{Sign Up Page}

It contains only the registration form. In the form must be entered first name, last name,
email address, password and password confirmation. Validation is performed in real time as the
user types, with error messages appearing beneath each field (if it's the case).

Email addresses are restricted to \texttt{@gmail.com} and \texttt{@yahoo.com} domains,
validated by splitting the input on the \texttt{@} character and checking the domain against an
allowed list. Passwords must be at least 8 characters long and contain at least one special
character, checked against a regular expression. The confirm password field is re-validated
whenever the password field changes, to immediately reflect any mismatch. All four validators
return a boolean and toggle CSS error classes on the input and its associated error span
(browser-native form validation is not used).

On form submission, \texttt{handleSignUp()} runs all four validators and only proceeds with
the API call if all pass. A POST request is sent to \texttt{/api/signup} with the form data
as JSON. On success, a toast notification instructs the user to check their email for the
verification link. The form fields are cleared but the user is not redirected, allowing them
to read the instruction. In the case of the server returning an error (example: because the
email is already registered) the error message from the server response is displayed in the
error toast.

When the user clicks the verification link in their email, it redirects back to
\texttt{sign\_up.html} with a \texttt{?verified=1} query parameter. The page detects this
parameter on load and displays a success toast confirming that the email has been verified and
that the user can now log in. The query parameter is then removed from the URL using
\texttt{history.replaceState} to keep the address bar clean.

The Sign Up page also embeds the login modal and forgot password modal, so users who already
have an account can log in directly from this page without navigating away.

\vspace{12pt}
\begin{figure}[H]
  \centering
  \includegraphics[width=0.75\textwidth]{signup_page}
  \caption{Figure 27 -- Sign Up page}
  \label{fig:signup}
\end{figure}
\vspace{12pt}

\subsubsection*{Login, Forgot Password, and Password Reset}

The login modal is available on every page, containing email and password fields. On successful
login, the header buttons update immediately to show Log Out and Admin Page, if it's the case.
If the logged-in account is the administrator, then the page redirects to
\texttt{admin.html} after a short delay. Login errors returned by the server such as unverified
accounts or wrong credentials are displayed inline beneath the password field rather than as a
generic toast.

The forget password modal can be accessed from a link that is present inside the login modal.
This takes an email id as input and POSTs to \texttt{/api/forgot-password}. The output will
always show a success message, even if the email id wasn't present in the database.

The password reset process is handled by a separate page, \texttt{reset\_password.html}. When
the user comes to the page using the reset link sent through the email, it takes the token from
the URL query string and passes it along with the reset request made to
\texttt{/api/reset-password}.

\vspace{12pt}
\begin{figure}[H]
  \centering
  \includegraphics[width=0.45\textwidth]{login_form}
  \caption{Figure 28a: Log In form}
  \label{fig:login}
\end{figure}
\vspace{12pt}

\begin{figure}[H]
  \centering
  \includegraphics[width=0.75\textwidth]{password_reset}
  \caption{Figure 28b: Password reset page}
  \label{fig:pwd-reset}
\end{figure}
\vspace{12pt}

\subsubsection*{Form Page}

Contains the form for contacting the website admin. It has fields for name, mail, subject and
message. Data is sent as JSON to \texttt{/api/contact}, then inserted into the database. A
success toast is displayed upon completion.

\subsubsection*{Predictions Page}

It fetches the live prediction data from \texttt{/api/predictions} endpoint on page load. The
predictions are in the form of day cards (built by \texttt{buildCard()} js function). Each card
has on it the day name, exact date, a sun or rain icon, predicted temperature and humidity. It
also has a bar for the safe or risk status. This bar shows the percentage related to the gas
concentration. The \texttt{formatDate()} function extracts date values in the format
YYYY-MM-DD from the API response dates. It formats values to be readable.

The page also contains informational text which explains when the predictions are generated and
details about how to interpret what's on the cards. The predictions are made to a total of up
to seven days.

\vspace{12pt}
\begin{figure}[H]
  \centering
  \includegraphics[width=0.85\textwidth]{predictions_page}
  \caption{Figure 29 -- Prediction page predictions view}
  \label{fig:predictions}
\end{figure}
\vspace{12pt}

\subsubsection*{Download Page}

Download page allows authenticated users to download the complete sensor dataset as a CSV file
and APK for the mobile application. For example, by clicking the download button, the frontend
sends a GET request to \texttt{/api/download-csv} (the case for CSV file) with the session
cookie included. In the case of the user not being logged in, the server returns a 401 response
and the frontend shows an error message prompting the user to log in first. When authenticated,
the server queries all rows from the \texttt{sensor\_data} table and writes them to an
in-memory buffer using Python's \texttt{csv} module. Afterwards it returns the buffer as a
downloadable file with the \texttt{Content-Disposition: attachment; filename=sensor\_data.csv}
header, so that the browser saves it directly as a file.

\vspace{12pt}
\begin{figure}[H]
  \centering
  \includegraphics[width=0.75\textwidth]{download_page}
  \caption{Figure 30 -- Zoomed view of the download page}
  \label{fig:download}
\end{figure}
\vspace{12pt}

\subsubsection*{Admin Page}

The Admin page is accessible only to the administrator account. It displays a table of all
registered users, showing their name, email address, verification status, registration date and
last login time. The list can be filtered by name, email, or registration date using a search
bar and date inputs that query the \texttt{/api/admin/users} endpoint. Maintenance mode can
also be toggled from here (works by calling \texttt{/api/admin/maintenance} with a POST
request). When maintenance mode is enabled, all non-admin visitors see a maintenance notice
instead of the website content.

In addition to the user management table and maintenance toggle previously described, the Admin
page also includes a Contact Messages panel from where all the information input by the users
on the form page can be seen. The information that can be seen is sender's name, email address,
subject, message text, and the timestamp at which the message was received. A sort order
dropdown allows the administrator to switch between newest-first and oldest-first ordering.

\vspace{12pt}
\begin{figure}[H]
  \centering
  \includegraphics[width=0.85\textwidth]{admin_page}
  \caption{Figure 31 -- Admin page view}
  \label{fig:admin}
\end{figure}
\vspace{12pt}

\section{PostgreSQL Database}
\label{sec:database}

\subsubsection*{Database Structure Overview}

The PostgreSQL database, named \texttt{website\_data}, is being hosted on AWS RDS and contains
eight tables. These can be grouped into three functional categories: sensor data storage, user
account management and site administration. It uses auto-incrementing integer primary keys
generated by sequences for most tables, with the exception of \texttt{site\_settings}, which
uses a text key as its primary key since it functions as a simple key-value configuration store.

\vspace{12pt}
\begin{figure}[H]
  \centering
  \includegraphics[width=0.85\textwidth]{er_diagram}
  \caption{Figure 32: Entity-relationship diagram of database}
  \label{fig:er-diagram}
\end{figure}
\vspace{12pt}

\subsubsection*{Sensor Data Table}

It stores the reading data and warnings. A new record is inserted every 30 minutes while the
station is active.

\vspace{12pt}
\begin{figure}[H]
  \centering
  \includegraphics[width=0.85\textwidth]{sensor_data_sql}
  \caption{Figure 33 -- sensor\_data table sql code}
  \label{fig:sensor-data-sql}
\end{figure}
\vspace{12pt}

The table stores three levels of detail for each MQ sensor: the raw 12-bit ADC integer value,
the computed voltage after voltage divider correction and the final PPM concentration. The raw
and intermediate values are kept as well, in case modifications to the formulas are made at a
later date. The timestamp fields sent by the station (year, month, day, time) are stored as
strings, matching the format in which the Raspberry Pi generates and sends them. The warning
column stores the full warning text in the case that any threshold is exceeded during that
cycle, or an empty string otherwise.

\subsubsection*{User Account Tables}

User account management is handled by two active tables: \texttt{users} and
\texttt{password\_reset\_tokens}.

\vspace{12pt}
\begin{figure}[H]
  \centering
  \includegraphics[width=0.85\textwidth]{users_sql}
  \caption{Figure 34 -- user table sql code}
  \label{fig:users-sql}
\end{figure}
\vspace{12pt}

The email column has a unique constraint, preventing duplicate registrations. Verified column
defaults to false at registration and is set to true only after the user clicks the
verification link sent by email. The \texttt{verify\_token} column stores the random token
sent in the verification email and is set to NULL after successful verification, invalidating
the link. The password column never stores the plain text password, only the bcrypt hash.
\texttt{Last\_login} timestamp is updated on every successful login.

The \texttt{password\_reset\_tokens} table manages the password reset flow.

\vspace{12pt}
\begin{figure}[H]
  \centering
  \includegraphics[width=0.85\textwidth]{prt_sql}
  \caption{Figure 35 -- password\_reset\_tokens table sql code}
  \label{fig:prt-sql}
\end{figure}
\vspace{12pt}

The \texttt{user\_id} column is a foreign key referencing \texttt{users(id)} with
\texttt{ON DELETE CASCADE}, meaning all reset tokens for a user are automatically deleted if
the user account is deleted. The token column has both a unique constraint and a B-tree index
(\texttt{idx\_prt\_token}). \texttt{User\_id} has its own index (\texttt{idx\_prt\_user\_id}).
These indexes ensure that token lookups and user-specific token deletions are fast operations.
When a new reset request is made, any existing token for that user is deleted before the new
one is inserted, ensuring only one valid reset token exists per user at any time.

\subsubsection*{Supporting Tables}

The \texttt{contact\_messages} stores the data, such as name, mail, subject and message and
exact sent time from the Contact page.

\vspace{12pt}
\begin{figure}[H]
  \centering
  \includegraphics[width=0.75\textwidth]{contact_messages_sql}
  \caption{Figure 36: contact\_messages table code}
  \label{fig:contact-sql}
\end{figure}
\vspace{12pt}

\texttt{site\_settings} table acts as a simple key-value store for Flask server.

\vspace{12pt}
\begin{figure}[H]
  \centering
  \includegraphics[width=0.75\textwidth]{site_settings_sql}
  \caption{Figure 37: site\_settings table code}
  \label{fig:settings-sql}
\end{figure}
\vspace{12pt}

The \texttt{aq\_predictions} table stores the 7-day forecasts generated by the prediction
script. The date column has a unique constraint, which enables the upsert behavior in the POST
endpoint. If the prediction script runs more than once on the same day, the existing row for
that date is updated. The \texttt{generated\_at} column stores the timestamp as a string in the
format produced by the Raspberry Pi (YYYY-MM-DD HH:MM:SS), while \texttt{received\_at} is set
automatically by the database server at the time of insertion or update.

\vspace{12pt}
\begin{figure}[H]
  \centering
  \includegraphics[width=0.85\textwidth]{aq_predictions_sql}
  \caption{Figure 38 -- prediction table sql code}
  \label{fig:pred-sql}
\end{figure}
\vspace{12pt}

Currently it stores a single entry with key \texttt{maintenance\_mode} and value true or false.
The key column is the primary key and new entries are inserted using PostgreSQL's
\texttt{ON CONFLICT (key) DO UPDATE} syntax, which acts as an upsert. So the maintenance mode
can be toggled without knowing whether the row already exists.

There are also three older tables, which were used for testing, at earlier stages of the
project: \texttt{form}, \texttt{log\_in} and \texttt{sign\_up}. These have been superseded by
\texttt{contact\_messages} and \texttt{users} respectively and are no longer written to by the
active application code.

\section{AWS Deployment}
\label{sec:aws}

\subsubsection*{EC2 Instance}

The Flask server runs on an Amazon EC2 instance located in the EU (Stockholm) region
(\texttt{eu-north-1}). The instance runs a Linux operating system and hosts the Flask
application on port 8080. EC2 instance serves a dual role: it handles all incoming HTTP
requests from both the Raspberry Pi (sensor data POST requests) and website visitors (page and
API requests). The Flask application is configured to bind to \texttt{0.0.0.0}, making it
accessible on all network interfaces, and the EC2 security group is configured to allow inbound
traffic on port 8080 from any IP address, so the website is publicly accessible.

\subsubsection*{RDS Instance and Connection}

The PostgreSQL database runs on an AWS RDS instance in the same \texttt{eu-north-1} region as
the EC2 instance. By hosting both in the same region network latency is kept to a minimum
between the Flask server and database. The RDS instance can be accessed at a fixed endpoint
hostname, so it lets the Flask application to connect to it by hostname, not IP address.

Flask server is connected to the database by using the \texttt{psycopg2} library through a
\texttt{SimpleConnectionPool}. It is configured to maintain between 1 and 10 open connections:

\vspace{12pt}
\begin{figure}[H]
  \centering
  \includegraphics[width=0.85\textwidth]{conn_pool_code}
  \caption{Figure 39: Data Base connection\_pool code}
  \label{fig:conn-pool}
\end{figure}
\vspace{12pt}

This is done using the \texttt{search\_path=public} parameter. By using it the table name
resolution will always refer to the public schema without having to specify it manually. API
routes that require database interaction invoke \texttt{get\_conn()} to get a connection from
the pool. \texttt{put\_conn()} is invoked at the end block to release the connection,
regardless of whether the interaction was successful or not. This is done to avoid leaking
connections in case of errors.

\vspace{12pt}
\begin{figure}[H]
  \centering
  \includegraphics[width=0.85\textwidth]{aws_deployment}
  \caption{Figure 40: AWS Deployment diagram}
  \label{fig:aws-deployment}
\end{figure}
\vspace{12pt}

\section{Android Mobile Application}
\label{sec:android}

The application is built using Flutter with the Dart programming language. It has four screens,
which are accessed through a bottom navigation bar: Home, Data, Predictions, and Contact.
Communication is done through the same Flask backend as the website using the same REST API
endpoints. No separate server side changes are needed.

The main entry point configures Google Fonts to use bundled fonts at build time
(\texttt{allowRuntimeFetching = false}) and launches the \texttt{AirNestApp} widget, which sets
up the Material.

\textbf{Home Screen}

The Home screen shows the same as the website page. Images are grouped as local resources in
the \texttt{assets/images} folder.

\textbf{Data Screen}

It structure and functionality are the same as the specific page on the website. Charts are
rendered using the \texttt{fl\_chart} library~\cite{flchart}, which renders line charts for
every parameter of the monitored sensors. Missing values in the dataset are treated the same
way as on the website when using Chart.js. Every chart has the same colors defined as on the
website.

\textbf{Predictions Screen}

Predictions are fetched from \texttt{/api/predictions}, same as for the website. Card styles
mimic styles on the predictions page of the website. DANGER cards have distinctive coloring.

\textbf{Contact Screen}

Contact screen, or ``Form page'' as it is on the website, has the same. When the form is
submitted, data is passed via JSON POST to the \texttt{/api/contact} endpoint of the Flask
server and saved to the \texttt{contact\_messages} table in the database. A success or error
message is shown when submitted.

\textbf{APK Distribution}

APK is built in release mode and can be downloaded from the Download page of the website
through the same process as the CSV download. It is distributed through a static file on the
EC2 instance along with the website files. Release build requires
\texttt{android:usesCleartextTraffic="true"} and \texttt{network\_security\_config.xml} to
make plain HTTP connection to the Flask server due to absence of HTTPS on the backend.
